\section{Methods}

\subsection{Data}

\paragraph{Primary cohort.} Two modalities from the same patient series were used.
Single-nucleus RNA sequencing (GSE308103) was generated with a fixed RNA-profiling probe
panel rather than a standard 3$'$ or 5$'$ chemistry, so it quantifies a defined set of
18{,}082 genes rather than the whole transcriptome; this is why housekeeping and
mitochondrial transcripts such as \emph{GAPDH}, the ribosomal protein genes and
\emph{MALAT1} are absent by design and are not used in any signature. The series contains
75 samples from 25 patients; 798{,}100 nuclei were called, of which 767{,}839 passed
quality control and 118{,}894 (15.48\%) were called as doublets, leaving
\textbf{413{,}697 nuclei} profiled over 18{,}069 genes for analysis.
Spatial transcriptomics (GSE307534) was generated on the Visium CytAssist platform in the
11\,mm format, 55\,$\mu$m spots on a $128\times112$ grid with a nominal spot-to-spot
distance of 100\,$\mu$m, giving 14{,}336 spots per section across \textbf{56 sections from
25 patients}; the platform implements the spatial-transcriptomics method of
\citet{stahl2016}. The two modalities share \textbf{23 patients}. The clinical sequence is
Normal $\rightarrow$ AAH $\rightarrow$ AIS $\rightarrow$ MIA $\rightarrow$ IAC; the spatial
annotations use the source study's token \texttt{LUAD} for the invasive stage, which we map
to IAC throughout. Section counts per stage are Normal 1, AAH 11, AIS 14, MIA 4, IAC 26
(Table~\ref{tab:cohort}).

\paragraph{External statistics.} Prognostic filtering used TCGA-LUAD
\citep{tcga2014luad} in its archived form: expression and clinical matrices only, giving
\textbf{483 patients and 177 deaths} (36.6\%; 381 patients and 143 events for disease-free
survival). The mutation archive in that distribution is a zero-byte file, so no
mutation-derived quantity (TMB, driver status) enters any analysis.

\paragraph{Perturbation resources.} Two libraries were queried. The compound library is
LINCS L1000 Level 5 (GSE70138; 118{,}050 signatures $\times$ 12{,}328 genes, 1{,}826
unique perturbagens across 41 cell lines) \citep{subramanian2017}; we rebuilt it locally
as described below rather than querying the CLUE service. The gene-level library is SCMG
\citep{pan2026scmg} (20{,}345 perturbations $\times$ 18{,}108 genes), used with the
authors' \texttt{CausalGenePredictor} and their reference cell-state manifold.

\paragraph{Annotation references.} Cell-type annotation was cross-checked against the
CellTypist \texttt{Human\_Lung\_Atlas} model \citep{dominguezconde2022celltypist}, itself
built on the HLCA \citep{sikkema2023}. Ligand analysis used the NicheNet human
ligand--receptor network and ligand--target matrix \citep{browaeys2020nichenet}.

\paragraph{Resources evaluated and not used.} Three external resources were assessed and
rejected, and we record them so that their absence is not read as an oversight. An external
normal-lung Visium series (GSE248082) was examined as a benign anchor for the copy-number
arm; its two usable sections failed a two-way health check (the held-out section scored
above threshold in both directions), so no external anchor is claimed. An external
lymph-node-metastasis series initially earmarked for the spatial cohort proved, on
inspection, to be breast rather than lung cancer and was removed. Three public single-cell
cohorts considered in the study's first design were moved out of scope when the paired
series became available, and were never analysed. No GWAS, eQTL, docking or structural data
entered this work.

\begin{table*}[!ht]
\centering
\begin{tabular*}{\textwidth}{@{\extracolsep{\fill}}l >{\raggedright\arraybackslash}p{172pt} >{\raggedright\arraybackslash}p{172pt}@{}}
\toprule
\tabheadfont Item & \tabheadfont Single-nucleus RNA-seq & \tabheadfont Spatial transcriptomics \\
\midrule
Accession & GSE308103 & GSE307534 \\
Platform & fixed RNA profiling & Visium CytAssist, 11\,mm \\
Feature space & 18{,}069 genes analysed & 18{,}085 probes \\
Samples & 75 & 56 sections \\
Patients & 25 & 25 \\
\addlinespace[3pt]
Units called & 798{,}100 nuclei & 14{,}336 spots per section \\
Retained after QC & 767{,}839 (96.2\%) & 95.03\% of spots \\
Doublet rate & 15.48\% & --- \\
Analysed & \textbf{413{,}697 nuclei} & all 56 sections \\
Grid and spacing & --- & $128\times112$; 100\,$\mu$m \\
\addlinespace[3pt]
Patients in both modalities & \multicolumn{2}{l}{23} \\
\bottomrule
\multicolumn{3}{@{}l}{\footnotesize\itshape Sections per stage: Normal 1 \quad AAH 11 \quad AIS 14 \quad MIA 4 \quad IAC 26 (source token \texttt{LUAD})} \\
\end{tabular*}
\caption{\textbf{Cohort and data sources.} Both modalities are public, come from the same
patient series, and share 23 patients. Both are probe-based assays rather than
whole-transcriptome chemistries, so each has a defined feature space. The invasive stage is
labelled \texttt{LUAD} in the source annotations and is mapped to IAC throughout. Normal is
represented by a single spatial section and MIA by four, so estimates for these two stages
should not be quoted on their own.}
\label{tab:cohort}
\end{table*}


% 参数表：与队列表相邻放在 Methods 前段，避免落到文末产生整页空白；
% 编号仍按正文首次引用顺序为 Table 6。
\begin{table*}[!ht]
\centering
\begin{tabularx}{\textwidth}{@{}>{\raggedright\arraybackslash}p{68pt}>{\raggedright\arraybackslash}p{78pt}>{\raggedright\arraybackslash}p{104pt}>{\centering\arraybackslash}p{20pt}X@{}}
\toprule
\tabheadfont Arm & \tabheadfont Parameter & \tabheadfont Value & \tabheadfont Src. & \tabheadfont Note \\
\midrule
\tabgroup{5}{Single-cell processing}
QC & retained nuclei & 399{,}579 & S & 15.48\% doublet rate; per-sample MAD \\
Clustering & L2 subtypes & 39 & S & endothelial and myeloid panels rebuilt \\
Integration & batch variable & sample ID & P & Harmony, rest at defaults \\
\addlinespace[3pt]
\tabgroup{5}{Spatial domains (BANKSY)}
Domain search & $\lambda$ & 0.2 & P & authors' recommendation for 55\,$\mu$m Visium \\
 & $k_{\text{geom}}$ & 18 & P & authors' recommendation \\
 & resolution & 0.5 & P & within the source study's range \\
Consensus & seeds & 5 & S & cross-seed agreement, per section \\
Archetypes & $K^{*}$ & 7 & S & stability 0.945 (0.814 at $K=8$) \\
\addlinespace[3pt]
\tabgroup{5}{Spatial deconvolution}
Reference & subtypes & 39 & S & a 6-lineage reference is also run \\
Epithelial call & rule & $\arg\max =$ epithelial & S & fixed before the cohort run \\
Spot mask & threshold & 500 counts / 200 genes / MT 0.15 & S & removes 4.97\% of spots \\
\addlinespace[3pt]
\tabgroup{5}{Copy-number inference}
Engine & tool & fastCNV & P & DOI 10.1186/s13073-026-01731-w \\
Chromosomes & filter & drop chrY, chrM & P & engine behaviour \\
Domain contrast & design & within-section pairing & S & removes depth by construction \\
\addlinespace[3pt]
\tabgroup{5}{Query signature}
Construction & axis & invasive $-$ precursor, median over patients & S & 16{,}104 genes \\
Survival filter & rule & disease-up needs HR$>$1; disease-down needs HR$<$1 & S & removes 51\% of genes \\
\addlinespace[3pt]
\tabgroup{5}{Perturbation querying}
Compound library & LINCS L1000 L5 & 118{,}050 $\times$ 12{,}328; 1{,}826 compounds & P & GSE70138 \\
 & scoring & CMap $\tau$; Spearman; Pearson & P & \texttt{signatureSearch} 1.12.0 \\
Gene library & SCMG & 20{,}345 perturbations & P & MIT-licensed \\
Gating & specificity & K562 control axis & S & removes 20 generalising entries \\
 & compartment & per-subtype axes (31) & S & top-100 counted per axis \\
Null models & permutations & 1{,}000 signature layer & T & \texttt{signatureSearch} default \\
Reproducibility & split-half & 4 partitions & S & reported individually \\
\bottomrule
\end{tabularx}
\caption{\textbf{Analysis parameters and their provenance.} Each parameter is labelled by
its source: \textbf{P}~= stated in a source paper or software publication, \textbf{T}~=
software default rather than a published recommendation, \textbf{S}~= fixed by us and
recorded in the analysis log. The distinction is drawn because most of the consequential
settings in this study are ours rather than the field's.}
\label{tab:params}
\end{table*}


\subsection{Single-cell processing}

After quality control and doublet removal 413,697 nuclei were retained. Contaminant and
low-quality clusters were then excluded, leaving 399,579 nuclei in the six mutually
exclusive lineage subsets that the atlas is built on. Expression was
normalised to counts per 10,000 with a $\log(1+x)$ transform (CP10K+log1p) throughout. The
atlas was built with Seurat \citep{hao2024seurat} using the regularised negative-binomial
variance-stabilising transform \citep{hafemeister2019sctransform} with \texttt{glmGamPoi}
\citep{ahlmanneltze2021glmgampoi} for parameter fitting, and integrated across samples with
Harmony \citep{korsunsky2019harmony}. Doublets were called with scDblFinder
\citep{germain2021scdblfinder}. An independent second caller was attempted but could not
serve as a cross-check: Scrublet \citep{wolock2019scrublet} reached a maximum score of
0.495 against its automatic threshold of 0.556 on these nuclei and therefore called zero
doublets in all 75 samples, a consequence of how sparse these profiles are rather than
evidence against the scDblFinder calls. The atlas was clustered per lineage, yielding 39
second-level
(L2) subtypes. Endothelial and myeloid marker panels were rebuilt before downstream
analysis: the published table provided six canonical genes per endothelial subtype, which
we expanded to as many as twenty genes per subtype across nine endothelial types; the
myeloid panel was expanded analogously, drawing on the human lung atlases
\citep{travaglini2020,habermann2020,vieirabraga2019,sikkema2023} and, for the epithelial
subtypes, on the lung adenocarcinoma atlas \citep{han2024,sinjab2021}. Automated annotation
against the HLCA reference used CellTypist \citep{dominguezconde2022celltypist}. All
reported subtype compositions use the expanded panels.

Ambient RNA was assessed in one precursor sample (P24) by three independent checks
(patient-level rank of the suspect programme, comparison against a matched negative-control
gene set, and an independent readout from a separate region of the same block). The sample
was retained with an annotation and all B/plasma-dependent conclusions are reported with
and without it.

\subsection{Spatial domain identification}

Spatial domains were identified with BANKSY \citep{singhal2024banksy} on each section
independently. The main
configuration was fixed at $\lambda=0.2$, $k_{\text{geom}}=18$ and resolution 0.5, chosen
because the authors recommend $\lambda\approx0.2$ for domain segmentation on 55\,$\mu$m
Visium and because higher $\lambda$ values, which the source study used, down-weight the
cell's own expression towards zero. Cross-seed consensus over five seeds produced 434
domain-level clusters across 56 sections; these were collapsed to archetypes by
profile-based clustering, giving $K^{*}=7$ at a subsampling stability of 0.945 (re-clustering on a random 80\% of
domains, 100 iterations); the cross-seed adjusted Rand index of the per-section domain
calls is a separate quantity and is lower (median 0.782). Stability
falls below 0.814 at $K=8$. A descriptive optimum exists at $K=6$ (0.966) but was not
adopted, since the selection rule took the largest $K$ passing the stability threshold.
Consensus stability was scored with the proportion-of-ambiguous-clustering statistic
\citep{senbabaoglu2014pac}.

A spatial-coherence criterion was also computed: per section, the mean coherence of the
partition was required to exceed the 95th percentile of a label-permutation null, with at
least 90\% of sections passing. This criterion is our own construction, not one recommended
by the source study or by the BANKSY authors, and it is not used to select the
configuration (see Discussion).

\subsection{Spatial deconvolution}

Spot composition was estimated with RCTD \citep{cable2022rctd} against a 39-subtype
reference, using the
full-section non-epithelial compartment as the reference for the copy-number arm (the
arm using each patient's own benign anchor covers only 10 of 25 cases). Spots were called
epithelial when the epithelial subtype had the highest weight; this criterion was fixed
before the cohort run. A spot mask was applied at thresholds of 500 counts, 200 genes and
0.15 mitochondrial fraction, removing 4.97\% of spots cohort-wide. This is higher than the
1.05\% reported by the source study; the difference is attributable to uneven per-section
depth rather than to the rule, and is recorded as a known discrepancy.

Two consequences of the mask are guarded explicitly. Per-section removal rates range from
0.02\% to 29.6\%, so a section in which a domain occupies little area may simply have had
more of itself masked; coverage is therefore reported per section and a small domain area is
never read as evidence that the niche is rare. And because the 56 sections come from 25
patients, every cross-section summary is clustered by patient, since two or three sections from
one patient are not independent samples.

\subsection{Depth matching}

Two independent estimators were used. In the first, spots were binned into per-slide deciles
of nUMI and compared only within bins. In the second, each gene was regressed on
$\log_{10}(\text{nUMI})$ per slide with a cubic term, and the domain-versus-rest ranking was
recomputed on the residuals. A gene was counted as surviving if it remained in the domain's
top 150 up-genes. The two estimators agree almost gene-for-gene.

\subsection{Spatial copy-number inference}

Copy-number scores were computed with fastCNV \citep{cabrejas2026fastcnv}. The engine
discards chrY and chrM, and its classification
output is a direction call rather than a malignancy determination; we use only the
continuous per-spot score. Domain-level comparisons pair each spot's score against the
median of the remaining spots in the same section, which removes depth by construction.
Significance was assessed per domain by a two-sided sign test over sections, with
Benjamini--Hochberg control where families of tests are reported
\citep{benjamini1995fdr}. Spatial neighbourhood statistics used Squidpy
\citep{palla2022squidpy}. Three further copy-number tools were evaluated and excluded
because each requires aligned reads rather than an expression matrix (Numbat
\citep{gao2023numbat}, HoneyBADGER \citep{fan2018honeybadger} and CaSpER
\citep{serinharmanci2020caspER}), as was SCEVAN \citep{defalco2023scevan}, which was
pre-registered but not run.

\subsection{Prognostic analysis}

Survival analysis used TCGA-LUAD \citep{tcga2014luad} (483 patients, 177 deaths;
disease-free survival: 381 patients, 143 events, censored at last follow-up), with Cox
proportional-hazards models fitted in R's \texttt{survival} package
\citep{therneau2000survival}. For each domain, the top-30 marker genes
were standardised across samples and averaged into a signature score. Univariate Cox models
were fitted both continuously (hazard ratio per standard deviation) and on a median split
(reference group: low). Hazard ratios for the adjusted models controlled for stage, age and
sex.

\subsection{Query signature construction}

For each patient, the disease axis was computed as the mean expression of invasive cells
minus that of precursor cells, and the median across patients was taken, giving an
18,069-gene vector, of which 16,104 are represented in the SCMG library. This axis was
filtered on independent survival evidence: a gene rising
in disease was retained only if its adjusted hazard ratio exceeded 1, and a gene falling in
disease only if it fell below 1; all others were set to zero. Of 18,069 genes, 8,799 were
retained (51\% removed overall; 44\% of the 15,792 genes with survival annotation).
A second, spatial query axis was constructed by comparing invasive and precursor spots
within nUMI deciles per section, merging deciles weighted by the smaller group size, and
removing immunoglobulin, mitochondrial and antibody-probe rows.

\subsection{A local compound-perturbation database}

The standard way to query a disease signature against LINCS is the CLUE service
(clue.io). Its scoring service was retired on 2026-01-31: the site still serves data
downloads, but it no longer accepts scoring jobs, and submission requests are refused
(HTTP 401). Rather than depend on a service that can no longer do the work, we rebuilt the
query resource locally, which has the further advantage that the exact library state is
fixed, versioned and re-runnable rather than whatever the remote service happens to serve
on a given day.

The database is a \texttt{signatureSearch}-format HDF5 file built directly from the LINCS
L1000 Level 5 Z-score matrix (GSE70138, release 2017-03-06). It stores the assay as a
genes $\times$ signatures array of 12{,}328 Entrez genes by 118{,}050 signatures, with
columns keyed by perturbagen, cell line and perturbation type; file size is 5.9\,GB.
Building it through R (\texttt{build\_custom\_db}) was abandoned because the required
compression and index handling failed and the per-column writer was prohibitively slow; the
file used here is written directly with \texttt{h5py}, which takes 26\,s. The database
itself is excluded from the repository, being 5.9\,GB and regenerable from the two tracked
build scripts. The \texttt{padj} slot is deliberately omitted because the scoring functions
used here do not read it.

Two access paths exist over the same underlying LINCS data, and it matters which analysis
used which. The disease-axis query, the spatial-axis query and the split-half stability
check were run through this local database using the official \texttt{signatureSearch}
package. The domain-signature query and all depth-matched sensitivity runs were run
directly on the source matrix through our own implementation of the CMap statistic. Both
paths are cross-validated against the official implementation (Table~\ref{tab:validation}).

\subsection{Slide registration}

The CytAssist images are the finer of the two available H\&E sources, but Space Ranger ships
no transform between them and the spot coordinate frame, so the mapping had to be measured.
We built a registration pipeline that places the high-resolution \texttt{hires} image into
the CytAssist pixel frame under a similarity transform with an optional horizontal
reflection,
$\text{cytassist} = s \cdot R(\theta) \cdot M \cdot \text{hires} + t$.

Three features make the search tractable rather than a blind optimisation. First, the scale
is not searched: it is computed analytically for each section from the shipped
scalefactors, $s_{\text{pred}} = (\text{fullres}\,\mu\text{m/px} \div \text{hires
scalefactor}) \div \text{CytAssist}\,\mu\text{m/px}$, with the CytAssist constant
4.6233\,$\mu$m/px; the search then only refines $\pm3\%$ around this value, and in the
final solution all 56 sections land within 1\% of it. Second, the rotation resolves to
essentially a binary choice: every section's fitted $\theta$ lies within 5$^\circ$ of
0$^\circ$ or 180$^\circ$, and the reflection is applied to 53 of 56 sections. Third, the
search is a coarse-to-fine pyramid over translation, obtained by phase correlation, scored
by a masked normalised cross-correlation of high-pass-filtered greyscale content (search
$\sigma = 12$, verification $\sigma = 4$).

Registration quality was established by three independent checks rather than one. The
adopted transform must score a fine-scale high-pass correlation of at least 0.10 and at
least three times the score of the competing reflection branch; this gate separates, by a
wide margin, the sections that the initial run failed on (the 52 passing sections have
correlation $\geq 0.19$ and separation $\geq 7.0$, while the four failures had separation
$\leq 1.4$ and were re-solved with a tissue-outline overlap metric). Second, the same
physical spot was cropped from both sources and compared directly, with the coordinate grid
propagated through the full transform; the fitted affine reproduces the grid to within
0.18\,px everywhere. Third, the residual was measured on the final solution as a translation field: the
residual translation is within one pixel on 46 of 56 sections, and three pixels or less on
the remainder, and where the field could be fitted it shows no residual rotation
($|\phi| < 0.1^\circ$) and no residual scale ($k$ within $\pm0.13\%$). At the CytAssist
scale of 4.6233\,$\mu$m per pixel this places even the largest residual well inside a spot
radius, so no spot is assigned to a neighbouring spot's tissue. All 56 sections pass. The
pipeline was used to align the same spots across image sources for the H\&E audit reported
in Results; it is retained in the repository.

\subsection{Perturbation scoring}

Compound scoring used the LINCS L1000 Level 5 library \citep{subramanian2017} (118,050
signatures $\times$ 12,328 genes; 1,826 compounds) through \texttt{signatureSearch}
\citep{duan2020signaturesearch} 1.12.0 with three schemes: the CMap weighted-KS statistic
($\tau$; the statistic originates in GSEA \citep{subramanian2005gsea}) on top-150 up/down
gene sets, and Spearman and Pearson correlation on the full numeric vector. Scores were
aggregated per compound by the median across cell lines. Gene-level scoring used the SCMG
library (20,345 gene perturbations)
with the authors' \texttt{CausalGenePredictor}
($\cos \times \text{perturbation sign} \times \text{gene-shift } z$), the authors' reference
manifold and their gene standard deviations.

\subsection{Positive control and numerical verification}

The pipeline was verified in two ways before interpretation. Numerically, our weighted-KS
implementation was compared against \texttt{signatureSearch} 1.12.0
(maximum absolute difference $6.6\times10^{-14}$; signs agreeing for 2000/2000
comparisons; the fast implementation used for the cross-check was \texttt{fgsea}
\citep{korotkevich2016fgsea}), and our vectorised causal score against the authors' Python
implementation
(Spearman 1.00000000 over 19,819 entries). Biologically, the official
epiblast-to-nascent-mesoderm benchmark was reproduced with the authors' own manifold and
standard deviations; the known regulators rank 1st (TBXT, of 8,450), 2nd (MSGN1), 3rd
(MIXL1), 7th (POU5F1), 9th (SNAI1), 11th (EOMES), 12th (NANOG) and 15th (EVX1).

\subsection{Candidate gating and compartment decomposition}

Gene-level candidates passed three additional gates. First, a specificity gate: each
perturbation was scored on a K562 control axis, and an entry whose effect there was as
strong as on the target axis was flagged as generalising and excluded; entries without
K562 data are reported as untested. Second, a compartment decomposition: the disease axis
was rebuilt independently for each of 31 subtypes (mean expression difference between
invasive and precursor cells within the subtype, median across patients) and every
perturbation was re-ranked within each axis; an entry was counted as top-100 for a subtype
at that rank. Third, a survival-direction check: a knockdown was kept only for genes whose
increased expression predicts worse outcome, and an overexpression only for the converse.

\subsection{Ligand analysis}

Ligand-level analysis used NicheNet \citep{browaeys2020nichenet} with the default human
ligand--receptor network and
ligand--target matrix. Fibroblasts were the receiver; the sender panel was the ten most
abundant non-fibroblast subtypes of the ECM/interstitial domain. The target set was the
top-200 up-genes of the data-derived disease axis, not a hand-written list. Candidate
ligands were required to be expressed above $\log(1+\text{CP10K})>0.1$ in at least one
sender and to be present in the ligand--receptor network.

\subsection{Reproducibility checks}

Split-half reproducibility used four random partitions of patients; each half was scored
against LINCS independently, and the Spearman correlation between the two rankings was
compared against a null formed by permuting one half's ranking (300 permutations).
Cross-modality agreement compared the single-cell-axis and spatial-axis compound rankings
by Spearman correlation and by top-$N$ overlap against a chance expectation of $N^2/118{,}050$.
Scale bars in spatial figures were calibrated per section from the median centre-to-centre
distance of adjacent in-tissue spots, which is nominally 100\,$\mu$m; this yields
1\,pixel $=0.2513\,\mu$m and a full-panel width of 10.9\,mm, consistent with the 11\,mm
capture format.

\subsection{Software}

Analyses used R 4.2.2 (Seurat, spacexr, BANKSY, Squidpy's R counterparts,
\texttt{signatureSearch} 1.12.0, survival), Python 3.8 (numpy, scipy, scanpy, h5py) and the
SCMG implementation \citep{pan2026scmg,xingjiepan_scmg}.
All figure code is in the accompanying repository. Table~\ref{tab:params} lists every
consequential parameter together with its provenance.


